% Administration: setting up a catalog, its settings, PostgreSQL, running it, the commands, one
% patient across databases, the Review tab, linking recordings, identifiers in recordings.
% Sources: docs/administration.md; medorama/src/config.rs (settings and defaults), main.rs,
% identity.rs, linkage.rs, review.rs, identify.rs, viewer.rs, server.rs (start_background).
\chapter{Administration}
\label{chap:administration}

This chapter is for the administrator who sets up a catalog and looks after it. A
\emph{catalog} is Medorama's own database of patients, encounters, clinical entries,
recordings, users and the audit log, together with its settings. The chapter describes where a
catalog lives, its settings, its database in SQLite or PostgreSQL, running the server, and the
decisions that keep one patient one patient across databases.

Other chapters cover the neighboring tasks. Users and Privacy Mode are described in
\cref{chap:users-privacy}, and the two ways of signing in, Medorama's passwords or the
institution's sign-in through a proxy, in \cref{chap:sign-in}. Deployment on servers and as a
service is described in \cref{chap:deployment}, and add-ons in \cref{chap:addons}. Every command
is described in full in \cref{chap:cli}, and every setting in \cref{chap:settings-reference}.
Bringing data in is the subject of \cref{part:data}: imported tables in \cref{chap:importing},
connected databases in \cref{chap:connecting}, and recordings in \cref{chap:recordings}.

\section{The rules}
\label{sec:administration-rules}

These rules hold for every deployment, and the code enforces each of them.
\begin{enumerate}
\item \textbf{Recordings are only read.} The catalog stores paths, never copies. It never
  writes, moves or deletes a recording, and never writes into a recording's folder, that is,
  the folder that holds a recording file, or anywhere inside a WFDB folder. An import refuses a
  recording that lies in such a position relative to the catalog's own folders (results,
  indexes, home).
\item \textbf{Removing affects the database only.} Removing a recording, an encounter, an import
  or a source deletes catalog entries. With \emph{purge}, the catalog's own outputs for a
  recording are deleted too, but only inside \code{results\_dir} and \code{index\_dir}. Every
  removal can be previewed first.
\item \textbf{Only administrators add or remove data and users.}
\item \textbf{A file comes in only with its patient.} Databases are the only way in: a
  recording, an ECG, an image or any other file enters as a row that names its patient (patient
  ID, MRN or encounter) and links the file. No file is imported on its own.
\item \textbf{Identifiers (PHI) live only in the catalog's database.} These are MRNs, names,
  dates of birth, encounter numbers, and what the recordings' own files say about the patient
  (\cref{sec:administration-identifiers}). They are sent, shown and searchable only for users with
  the PHI permission who switched Privacy Mode off, and those users' searches and views are
  written to the audit log. The viewers, the job descriptions given to actions, the results
  folders and the indexes carry pseudonyms only.
\item \textbf{Outputs go to configured folders.} Everything the catalog writes goes to
  \code{results\_dir}, \code{index\_dir} or its home folder, all chosen when the catalog is
  created.
\end{enumerate}

\section{Setting up a catalog}
\label{sec:administration-catalog}

A catalog needs two places, chosen when it is created (\cref{tab:administration-places}).

\begin{xltabular}{\linewidth}{@{}P{2.4cm}P{4.6cm}L@{}}
\caption{The two places of a catalog.}
\label{tab:administration-places}\\
\toprule
\tabhead Place & \tabhead What goes there & \tabhead Advice \\
\midrule
\endfirsthead
\toprule
\tabhead Place & \tabhead What goes there & \tabhead Advice \\
\midrule
\endhead
\bottomrule
\endfoot
Home folder & \file{medorama.toml} (the settings), \file{medorama.sqlite} (the database),
  \file{actions/} (plug-ins), \file{formats/} (import formats), \file{viewer/} (the settings and
  log of the viewer the catalog runs), \file{secrets/} (connected databases' passwords) &
  On a \textbf{local disk}: SQLite's locking is unreliable on network file systems such as NFS
  or SMB home folders. Private to the account that runs the catalog, and encrypted: the
  database holds identifiers. Default: \file{\textasciitilde/.local/share/medorama} (Windows:
  \code{\%LOCALAPPDATA\%\textbackslash medorama}). \\
Results folder & \file{index/} (the recordings' indexes) and \file{actions/} (the results of
  actions) & Room for the indexes, roughly 1 to 5\,\% of the recordings' size; never a folder
  that holds recordings. \\
\end{xltabular}

\cmd{medorama init} creates the catalog:
\begin{lstlisting}[style=shell]
medorama init --results /data/medorama-results                     # the default home folder
medorama init --results /data/medorama-results --home ~/medorama   # or one of your own
\end{lstlisting}
The command asks twice for the password of the first user, \code{admin} (at least 10
characters), who is an administrator with the PHI permission. It writes a commented
\file{medorama.toml}, creates the database, makes the home folder readable by its owner only, and
remembers the new catalog as the default for later commands. When users will sign in through the
institution's proxy, name the first administrator as the proxy will name that person, with
\opt{-{}-admin} (\cref{sec:sign-in-accounts}). Its options (\opt{-{}-index-dir}, \opt{-{}-admin},
\opt{-{}-timezone}, \opt{-{}-port}, \opt{-{}-password-stdin}, \opt{-{}-database-url},
\opt{-{}-database-password-file}, \opt{-{}-no-demo}) are listed in \cref{chap:cli}. A catalog
starts empty unless the installer was given \opt{-{}-demo}: the demo, its use and its removal are
described in \cref{chap:quickstart}; never install it where real patients are.

\subsection{Which catalog a command uses}
\label{sec:administration-which}

Every command works on one catalog: the first of these that applies.
\begin{enumerate}
\item the global option \opt{-{}-home DIR}, given before or after the command;
\item the environment variable \code{MEDORAMA\_HOME};
\item the catalog that \cmd{init} created last, or the one that \cmd{medorama -{}-home DIR use}
  chose, remembered in \file{\textasciitilde/.config/medorama/home} (Windows:
  \code{\%APPDATA\%\textbackslash medorama\textbackslash home}), if that folder still holds a
  \file{medorama.toml};
\item the built-in home folder, \file{\textasciitilde/.local/share/medorama} (Windows:
  \code{\%LOCALAPPDATA\%\textbackslash medorama}).
\end{enumerate}
\cmd{medorama serve} prints the catalog it uses, with its address and its viewer (\code{viewer:
  ... (started on first use, in the VitalStream tab)}). A command run where there is no catalog says so
(\code{no catalog in ...}) and explains how to create one or name another.

\section{The settings: medorama.toml}
\label{sec:administration-settings}

The settings are in \file{<home>/medorama.toml}, which \cmd{init} writes with a comment above
each setting. Relative paths are relative to the home folder; \code{\textasciitilde} stands for
the home folder of the account. The server reads the file when it starts, so restart it after a
change. A key that the program does not know is an error (\code{unknown settings: ...}), so that
a mistyped setting is never silently ignored, and so is a value of the wrong type
(\code{bad value for ...}). \Cref{tab:administration-settings} lists every setting with its default; the types and
the details are in \cref{chap:settings-reference}.

\begin{xltabular}{\linewidth}{@{}P{3.7cm}P{3.3cm}L@{}}
\caption{The settings of \file{medorama.toml}, with their defaults.}
\label{tab:administration-settings}\\
\toprule
\tabhead Setting & \tabhead Default & \tabhead Meaning \\
\midrule
\endfirsthead
\toprule
\tabhead Setting & \tabhead Default & \tabhead Meaning \\
\midrule
\endhead
\bottomrule
\endfoot
\code{results\_dir} & (set by \cmd{init}; required) & Where action results go:
  \file{actions/<action>/r<recording>/j<job>/} \\
\code{index\_dir} & \file{<results\_dir>/index} & The recordings' indexes, one folder per index
  key \\
\code{adapters\_dir} & \file{\textasciitilde/.local/share/medorama/adapters} (Windows:
  \code{\%LOCALAPPDATA\%\textbackslash medorama\textbackslash adapters}) & The adapter folder
  (\cref{chap:addons}) \\
\code{database} & \file{medorama.sqlite} & The SQLite database file; it holds identifiers \\
\code{database\_url} & (empty: SQLite) & A PostgreSQL database instead,
  \code{postgres://user@host/db?sslmode=require}, without its password;
  \code{MEDORAMA\_DATABASE\_URL} overrides it \\
\code{database\_password\_file} & (none) & The file that holds that database's password;
  \code{MEDORAMA\_DATABASE\_PASSWORD} overrides it \\
\code{database\_ca\_file} & (none) & Certificate authorities (PEM) trusted for the database
  server's certificate, besides the system's \\
\code{host}, \code{port} & \code{127.0.0.1}, 8090 & Where the web interface and the API
  listen \\
\code{workers} & 2 & Jobs run at the same time (a value below 1 is read as 1) \\
\code{auto\_index} & \code{true} & While the server runs, queue the indexing of new and changed
  recordings \\
\code{check\_interval\_s} & 900 & Re-check every recording this often, in seconds (0: never) \\
\code{timezone} & \code{America/New\_York} & The time zone of times written without an offset,
  and of the times shown \\
\code{viewer} & \code{auto} & The VitalStream viewer the catalog runs: \code{auto}, \code{cpp},
  \code{rs}, \code{py}, \code{none}, or a program's path (\cref{sec:administration-viewer}) \\
\code{viewer\_url} & (empty) & Instead, a viewer running elsewhere, opened in a new browser
  tab \\
\code{qc\_verbose} & \code{false} & The viewer the catalog runs also marks (dashed outline) the
  data a reader repaired, such as data blocks a file stored out of order; breaks are shaded
  either way \\
\code{data\_roots} & \code{[]} (anywhere) & The folders files may be imported from \\
\code{actions\_dirs} & \code{["actions"]} & The folders of action plug-ins
  (\cref{chap:actions}) \\
\code{formats\_dirs} & \code{["formats"]} & The folders of import formats described in TOML
  (\cref{chap:importing}) \\
\code{session\_hours} & 12 & How long a sign-in lasts, in hours (a value below 1 is read as
  1) \\
\code{session\_idle\_minutes} & 60 & A sign-in unused this long ends (0: only
  \code{session\_hours} ends it); the page's own checks for new data do not count as use \\
\code{secure\_cookies} & \code{false} & \code{true} when the catalog is served over HTTPS \\
\code{sign\_in} & \code{"password"} & How users sign in: \code{"password"} (Medorama's accounts)
  or \code{"proxy"} (the institution's authenticating proxy; \cref{chap:sign-in}) \\
\code{proxy\_user\_header} & \code{"X-Remote-User"} & Proxy mode: the header that names the
  user \\
\code{proxy\_addresses} & \code{["127.0.0.1", "::1"]} & Proxy mode: the IP addresses the proxy
  connects from; the header is believed from these only \\
\code{proxy\_logout\_url} & (empty) & Proxy mode: where \ui{Sign out} leads, the proxy's own
  sign-out \\
\code{key\_file} & (none) & The database key, for builds with encryption
  (\cref{chap:deployment}) \\
\code{audit\_export} & (none) & A file the audit log's sealed entries are appended to, on
  storage the database's administrators cannot change (\cref{chap:users-privacy}) \\
\end{xltabular}

A new catalog's file, as \cmd{init} writes it (the comments shortened here):
\begin{lstlisting}[language=TOML]
results_dir = "/data/medorama-results"
# index_dir = "/data/medorama-results/index"
adapters_dir = "/home/medorama/.local/share/medorama/adapters"
# database = "medorama.sqlite"
host = "127.0.0.1"
port = 8090
workers = 2
auto_index = true
check_interval_s = 900
timezone = "America/New_York"
viewer = "auto"
viewer_url = ""
qc_verbose = false
data_roots = []
actions_dirs = ["actions"]
formats_dirs = ["formats"]
session_hours = 12
session_idle_minutes = 60
secure_cookies = false
# sign_in = "proxy"
# proxy_user_header = "X-Remote-User"
# proxy_addresses = ["127.0.0.1", "::1"]
# proxy_logout_url = "/oauth2/sign_out"
# audit_export = "/var/log/medorama/audit.jsonl"
\end{lstlisting}

\begin{tip}
Set \code{data\_roots} before the first import, to the folders that hold the data and nothing
more: files outside them are refused at import, and the folders a user may browse are limited
to them. \cmd{medorama import-clif} is the exception: unlike the page and the API, the command
does not check its folder against \code{data\_roots}.
\end{tip}

Medorama has no setting for the limits on a WFDB folder (5000 header files, 20 folder levels,
10\,s to count them): it indexes a recording's folder within those defaults. They are settings of
the VitalStream viewers (\code{max\_records}, \code{max\_depth}, \code{max\_scan\_s};
\cref{chap:settings-reference}).

\section{The catalog in PostgreSQL}
\label{sec:administration-postgres}

SQLite, the default, suits one server: there is nothing to install, and the file can be
encrypted. The catalog goes into \textbf{PostgreSQL} (14 or later) to run several Medorama
servers on one catalog, for many users or large catalogs, or where the institution runs
PostgreSQL already (on its own servers, Amazon RDS or Aurora, Azure Database for PostgreSQL).
The same program does both, and every test runs on both.

\begin{enumerate}
\item \textbf{An empty database}, with an account for Medorama alone, owner of it:
\begin{lstlisting}[language=SQL]
CREATE ROLE medorama LOGIN PASSWORD '...';
CREATE DATABASE medorama OWNER medorama;
\c medorama
CREATE EXTENSION IF NOT EXISTS pg_trgm;   -- the text search's index (needs a superuser once)
\end{lstlisting}
\item \textbf{The password} in a file that only the catalog's account can read
  (\cmd{chmod 600}), or in \code{MEDORAMA\_DATABASE\_PASSWORD}; never in the URL: a
  \code{database\_url} that holds a password is refused.
\item \textbf{Create the catalog}; its settings and folders stay in the home folder:
\begin{lstlisting}[style=shell]
medorama init --results /data/medorama-results \
    --database-url "postgres://medorama@db.example.org/medorama" \
    --database-password-file /etc/medorama/db-password
\end{lstlisting}
  A database server on another host must offer TLS: Medorama requires it (unless the URL says
  \code{sslmode=disable}) and verifies the server's certificate and name.
  \code{database\_ca\_file} adds an authority (a private CA, Amazon RDS's bundle).
\item \textbf{Several servers:} give each the same \file{medorama.toml} (the same
  \code{database\_url} and \code{data\_roots}, the results and index folders on shared storage),
  behind a reverse proxy (\cref{chap:deployment}). Any server answers any request; jobs are
  taken one at a time from the shared queue; schema updates, taking a job, the automatic
  indexing, imports and every change to patients and encounters (an import's chunk, a merge, a
  connected database's persons linked) are done by one server at a time (advisory locks), so
  that one MRN never names two patients. A new SQLite catalog opened by several processes at
  once is created by one of them.
\item \textbf{Connections:} each server opens at most 32 connections to the database (a request
  waits for one beyond that), checks a connection that was idle before using it again, and,
  unless the URL sets them, gives each connection 10 seconds to connect, TCP keepalives, and 5
  minutes for a statement (\code{?connect\_timeout=\dots}, \code{keepalives=0},
  \code{options=-c statement\_timeout=\dots} in the URL change them). The URL holds no password,
  neither before the host nor as a parameter: Medorama refuses one that does.
\end{enumerate}

\textbf{A SQLite catalog moves to PostgreSQL} with every table and row, in one transaction (each
table read a range of rows at a time, never whole), and each database's table of numbers
(\code{dbcol\_<import>}) is built again in PostgreSQL:
\begin{lstlisting}[style=shell]
medorama copy-database --to "postgres://medorama@db.example.org/medorama" \
    --password-file /etc/medorama/db-password
\end{lstlisting}
It prints the lines to add to \file{medorama.toml}. The SQLite file is left as it was: keep it as
a backup, then delete it securely.

\section{Running it}
\label{sec:administration-running}

\begin{lstlisting}[style=shell]
medorama serve
\end{lstlisting}
Open \url{http://127.0.0.1:8090} and sign in. \key{Ctrl}+\key{C} stops the server (as does the
\code{SIGTERM} of a service manager's stop): its plug-ins are stopped (TERM to each one's process
group, KILL 5 seconds later), the jobs it was running go back to the queue, and the viewer it
started stops with it; \cmd{medorama work} stops in the same way. At start-up, an import that a
stop or a crash interrupted is removed (\cref{sec:importing-rule}). While it runs, in the
background:
\begin{itemize}
\item \textbf{workers} (\code{workers}) run the queued jobs: indexing, derived series, actions;
\item \textbf{automatic indexing} (\code{auto\_index}) queues new and changed recordings, every
  15 seconds;
\item every recording is \textbf{re-checked} every \code{check\_interval\_s} seconds: missing
  recordings, changed ones, lost indexes;
\item the audit log is \textbf{sealed} every minute (\cref{chap:users-privacy}).
\end{itemize}
Without the server, the commands do the same work: \cmd{medorama run index -{}-all} queues the
indexing of every recording, and \cmd{medorama work} runs the queue in that process until it is
empty. \cmd{medorama serve -{}-no-workers} serves the page only, while another process runs
\cmd{work}.

When the server listens on an address other than \code{127.0.0.1}, \code{localhost} or
\code{::1} without \code{secure\_cookies = true}, it prints a warning: the page would be reachable
from other machines without HTTPS. \Cref{chap:deployment} describes the ways to reach a server
from elsewhere.

\subsection{The waveform viewer}
\label{sec:administration-viewer}

Medorama runs one VitalStream viewer for its users, and shows it in its \ui{VitalStream} tab.
The viewer programs are the standard ones, unchanged; the setting \code{viewer} chooses which one
(\cref{tab:administration-viewer}).

\begin{xltabular}{\linewidth}{@{}P{3.4cm}L@{}}
\caption{The setting \code{viewer}.}
\label{tab:administration-viewer}\\
\toprule
\tabhead Setting & \tabhead Viewer \\
\midrule
\endfirsthead
\toprule
\tabhead Setting & \tabhead Viewer \\
\midrule
\endhead
\bottomrule
\endfoot
\code{auto} (default) & The one that \cmd{tools/install-medorama.sh} installed (recorded in
  \file{medorama.viewer} next to Medorama's program), else the first found of
  \cmd{vitalstream-cpp}, \cmd{vitalstream-rs}, \cmd{vitalstream-py} \\
\code{cpp}, \code{rs}, \code{py} & That viewer \\
A path & That program \\
\code{none} & No viewer: \ui{View} and the \ui{VitalStream} tab are grayed out, with the
  reason \\
(\code{viewer\_url} set) & The catalog runs none, and \ui{View} opens that viewer in a new
  browser tab \\
\end{xltabular}

\paragraph{How it runs.} At the first \ui{View}, or the first visit of the tab, the catalog
starts the viewer as a process of its own on \code{127.0.0.1} and a free port. It restarts the
viewer if its process ends, or if a request cannot reach it and it then fails a health check, and
stops it when the catalog stops. A request to a running viewer waits for no other request and
costs no check, so a viewer busy with a large recording is never taken for a dead one. A viewer left running by an earlier run
that ended otherwise is stopped and a new one started, so that the viewer always has the
catalog's current settings and add-ons.

The viewer's settings, \file{<home>/viewer/vitalstream.toml}, are rewritten at each start. They
give it the catalog's \code{index\_dir} (every viewer writes the same index), its
\code{adapters\_dir}, its time zone, its \code{data\_roots} and \code{qc\_verbose}, and
\code{require\_password = true}, so that a viewer built without password support refuses to run
under the catalog. The viewer's password file, \file{<home>/viewer/auth.json}, has a signing key
that only the catalog knows. The catalog signs every request it passes on, users never see the
viewer's sign-in, and the viewer's port is of no use to anyone else on the host. The viewer's
output goes to \file{<home>/viewer/viewer.log}.

\paragraph{The gateway.} The catalog passes the viewer's requests, at \file{/viewer/}, only after
its own sign-in, and only the requests it lists (default deny), each checked against the user's
rights. Administrators open recordings inside \code{data\_roots}; other users open only the
recordings of the catalog. Every recording opened is written to the audit log (\code{view}). In
Privacy Mode, recordings are opened by token, folders are not browsed, the gateway moves times
both ways, and the texts of EDF+ annotations are treated like notes
(\cref{sec:protecting-edf}). The gateway
answers 413 to a request body over 1\,MiB, and 415 to a \code{POST} whose body is not sent as
\code{application/json} (a body that is not valid JSON: 400); a \code{POST} must also carry
\code{X-Medorama: 1}, as every changing request to the catalog does. In Privacy Mode, the
viewer's error texts, which name files, are replaced by a fixed sentence. The rules, request by
request, are in \cref{chap:viewer-api}.

\paragraph{A viewer run by hand} (for oneself, on another port) can share the catalog's indexes:
set its \code{cache\_dir} to the catalog's \code{index\_dir} and its \code{adapters\_dir} to the
catalog's.

\section{The commands}
\label{sec:administration-commands}

\opt{-{}-home DIR} is a global option: every command takes it, before or after the command's
name. \cmd{medorama <command> -{}-help} shows its options.
\Cref{tab:administration-commands} is an overview; \cref{chap:cli} gives each command and option
in full.

\begin{xltabular}{\linewidth}{@{}P{7.4cm}L@{}}
\caption{The commands of \cmd{medorama}, in brief.}
\label{tab:administration-commands}\\
\toprule
\tabhead Command & \tabhead Does \\
\midrule
\endfirsthead
\toprule
\tabhead Command & \tabhead Does \\
\midrule
\endhead
\bottomrule
\endfoot
\cmd{init -{}-results DIR ...} & Create a catalog \\
\cmd{-{}-home DIR use} & Make that catalog the default for later commands \\
\cmd{serve [-{}-host H] [-{}-port N] [-{}-no-workers]} & The web interface, the API, the
  background work \\
\cmd{import-table FILE... [-{}-name N] [-{}-source S] [-{}-category C] [-{}-mapping F]
  [-{}-format F] [-{}-dry-run]} & Import a database: CSV, Parquet or a spreadsheet
  (\cref{chap:importing}) \\
\cmd{import-clif FOLDER [-{}-source S] [-{}-dry-run]} & A CLIF folder: each table a database
  (\cref{chap:clif}). Unlike the page and the API, the command does not check the folder against
  \code{data\_roots} \\
\cmd{formats [-{}-check FILE]} & The import formats; check a format file \\
\cmd{imports list}, \cmd{imports remove ID [-{}-purge-results] [-{}-dry-run]} & The imports;
  remove one with what it added \\
\cmd{sources list}, \cmd{sources remove NAME [-{}-purge-results] [-{}-dry-run]} & The sources;
  remove everything from one \\
\cmd{connections add NAME -{}-kind omop|sql -{}-url URL ...}, \cmd{list}, \cmd{check NAME},
  \cmd{sync NAME}, \cmd{remove NAME} & Connected databases (\cref{chap:connecting}) \\
\cmd{demo add [DIR]}, \cmd{demo remove [-{}-dry-run]} & Load the demo (the installer's, or a
  folder made by \file{tools/make-demo-data.py}); remove it with everything it added
  (\cref{chap:quickstart}) \\
\cmd{addons list}, \cmd{addons install PROGRAM}, \cmd{addons remove NAME} & Add-ons
  (\cref{chap:addons}) \\
\cmd{list [-{}-q SEARCH] [-{}-unit U] [-{}-status S] [-{}-limit N] [-{}-phi] [-{}-json]} &
  encounters; \opt{-{}-limit} default 100; identifiers only with \opt{-{}-phi} \\
\cmd{recordings [-{}-q SEARCH] [-{}-unlinked] [-{}-limit N] [-{}-json]} & Recordings;
  \opt{-{}-limit} default 100 \\
\cmd{remove ID|PATH... [-{}-search Q] [-{}-purge-results] [-{}-dry-run]} & Remove recordings from
  the catalog (the files are untouched) \\
\cmd{remove-encounters ID...|-{}-search Q|-{}-all [-{}-dry-run]} & Remove encounters with their
  bed stays and items (the recordings stay) \\
\cmd{check} & Re-check the recordings: missing, changed, lost indexes \\
\cmd{actions} & The installed actions \\
\cmd{run ACTION... -{}-all|-{}-encounter ID...|-{}-recording ID...|-{}-search TEXT|-{}-cohort N
  [-{}-force] [-{}-param A.K=V] [-{}-priority N]} & Queue actions \\
\cmd{work [-{}-workers N]} & Run the queue in this process until it is empty \\
\cmd{jobs [-{}-state S] [-{}-limit N] [-{}-json]}, \cmd{cancel -{}-job N|-{}-run N|-{}-all} & The
  queue (\opt{-{}-limit} default 50); \cmd{cancel -{}-job N -{}-run M} cancels job N only if it
  belongs to run M \\
\cmd{user add NAME -{}-role R [-{}-phi] [-{}-display-name D] [-{}-password-stdin|-{}-no-password]},
  \cmd{user list}, \cmd{user set NAME [-{}-role R] [-{}-phi true|false] [-{}-disabled true|false]
  [-{}-display-name D] [-{}-password|-{}-password-stdin]}, \cmd{user remove NAME} & Users
  (\cref{chap:users-privacy}); \cmd{user set}'s \opt{-{}-phi} and \opt{-{}-disabled} take a value
  (\opt{-{}-disabled true}) \\
\cmd{audit [-{}-limit N]}, \cmd{audit -{}-verify}, \cmd{audit -{}-seal},
  \cmd{audit -{}-seal -{}-verify} & The audit log (\opt{-{}-limit} default 50); whether its seals
  hold; seal it now; seal, then check \\
\cmd{copy-database -{}-to URL [-{}-password-file F] [-{}-ca-file F]} & Copy a SQLite catalog into
  an empty PostgreSQL one (the password file, else \code{MEDORAMA\_DATABASE\_PASSWORD};
  \opt{-{}-ca-file}: certificate authorities, PEM, besides the system's) \\
\end{xltabular}

\paragraph{Rights.} The commands act directly on the database, with administrator rights and
the PHI permission, as the operating-system user who can read the home folder; the audit log
records them as \code{local:<user>}. \cmd{list} shows identifiers only with \opt{-{}-phi}.
Whoever can read the home folder can read the database anyway, so run the catalog under a
dedicated account, and let people use the web page.

\section{One patient across databases}
\label{sec:administration-identity}

The catalog keeps one patient one patient across databases, in the way of a master patient
index: what each database said is kept, the catalog's view is computed from it, and decisions are
kept apart from the data.

\begin{itemize}
\item \textbf{Matching patients:} by \code{patient\_id} (or an ID merged into another one), else
  by an MRN of the same authority (\code{mrn\_authority}; ignoring case and surrounding spaces).
  A row with an MRN but no \code{patient\_id} gets a pseudonymous ID from the catalog
  (\code{GEN-000001}, \code{GEN-000002}, \dots), the same for that MRN after any removal and
  re-import, and never given to another.
\item \textbf{Values:} each field shows the newest batch's value; removing that batch brings back
  the previous one. A patient may have MRNs of several authorities; \code{mrn:} searches find any
  of them.
\item \textbf{A row is held back} when it contradicts the catalog: kept, with a warning, for the
  Review tab (\cref{sec:administration-review}), and counted as a conflict in the import's
  preview (\cref{tab:administration-held}).
\end{itemize}

\begin{xltabular}{\linewidth}{@{}P{6.6cm}L@{}}
\caption{Rows held back, and what applying them does.}
\label{tab:administration-held}\\
\toprule
\tabhead Problem & \tabhead Applying it \\
\midrule
\endfirsthead
\toprule
\tabhead Problem & \tabhead Applying it \\
\midrule
\endhead
\bottomrule
\endfoot
Its \code{patient\_id} and its MRN belong to two different patients & Merges the MRN's patient
  into the ID's \\
Its MRN belongs to a patient with another \code{patient\_id} & Records the two IDs as one
  person \\
The patient already has another MRN of the same authority & Adds the MRN \\
Its encounter belongs to another patient & Moves the encounter, with its items \\
\end{xltabular}

The catalog never re-assigns a patient on its own on the strength of an MRN. \textbf{A merge}
moves the patient's encounters, MRNs, values, items and sources to the other patient; later rows
that name the merged ID go to the other patient, and the merged patient's pseudonym stays an
alias of the survivor. A merge cannot be undone: removing and re-importing the batches does not
split it either.

\section{The Review tab}
\label{sec:administration-review}

The \ui{Review} tab, for administrators, gathers what needs a person's decision
(\cref{tab:administration-review}). Every decision is kept, outlives re-imports, and is written
to the audit log (\code{apply\_held\_row}, \code{dismiss\_held\_row}, \code{merge\_patients},
\code{distinct\_patients}, \code{link\_decision}).

\begin{xltabular}{\linewidth}{@{}P{3.2cm}P{5.6cm}L@{}}
\caption{The lists of the Review tab.}
\label{tab:administration-review}\\
\toprule
\tabhead List & \tabhead Shows & \tabhead Decisions \\
\midrule
\endfirsthead
\toprule
\tabhead List & \tabhead Shows & \tabhead Decisions \\
\midrule
\endhead
\bottomrule
\endfoot
\ui{Held-back rows} & Rows in conflict with the catalog (\cref{tab:administration-held});
  without the PHI permission, their columns without identifiers & Apply (with the fix shown) or
  dismiss \\
\ui{Possible duplicate patients} (PHI users) & The same name and date of birth; first and last
  names swapped; day and month of birth swapped; the same last name, first initial and date of
  birth & Merge into either one; different people (never suggested again) \\
\ui{Recordings whose files disagree} (PHI users) & Automatic links where the recording's own
  file names another patient (\cref{sec:administration-identifiers}) & Keep the link (it becomes
  manual, and is not shown again); unlink (never linked automatically again) \\
\ui{Items waiting} & Items that no database names yet, grouped by the keys they wait with &
  none: they are matched when a database names their patient or encounter \\
\end{xltabular}

\section{Linking recordings}
\label{sec:administration-linking}

A recording is linked to an encounter after every import and every indexing:
\begin{enumerate}
\item \textbf{by a database row} that names its file and the encounter (the \code{file}
  column), or by a bed stay's \code{recording\_file};
\item \textbf{by bed and time:} the unit and bed of its name (bedside-monitor systems name files
  \code{UNIT\_BED-<Unix time>}) equal a bed stay's \code{care\_unit}, an underscore and its
  \code{bed} (ignoring case), and its time span overlaps the stay: before indexing, the time in
  its name; after, its real start and end.
\end{enumerate}
\textbf{Decisions made manually} win over both, and outlive removing and re-importing the
databases: \textbf{link} (the recording gets exactly its manual links) and \textbf{never}
(\ui{Unlink}: never linked to that encounter automatically). The recordings that match nothing
are listed under \ui{Import} $\rightarrow$ \ui{Recordings} $\rightarrow$ \ui{Not linked to an
encounter}, and by \cmd{medorama recordings -{}-unlinked}. When recordings are not linked, check
the databases' times first: a time zone that differs from the setting \code{timezone} moves every
bed stay.

\section{Identifiers in recordings}
\label{sec:administration-identifiers}

Some recording files store the patient's identifiers themselves. Bedside monitors write them
into their files, and an EDF or BDF header has a patient field and a recording field. The
built-in action \code{identify} reads these identifiers, once per recording and again when the
file changes, so that an administrator can check that each file belongs to the patient it is
linked to. WFDB records hold none.
\begin{itemize}
\item \textbf{EDF, EDF+, BDF and BDF+ files} are read by Medorama itself (\code{from\_edf} in
  \file{medorama/src/identify.rs}). EDF+ structures the patient field as
  \code{code sex birthdate name} and the recording field as
  \code{Startdate dd-MMM-yyyy admincode technician equipment}; \code{identify} keeps the code
  (as the patient ID), the sex, the date of birth, the name and the admin code (as the encounter
  ID), and leaves out the parts written \code{X} (unknown). A plain EDF or BDF patient field has
  no structure and is kept whole, as the patient ID. The engine and the viewers never put these
  fields into an index or any answer (\cref{sec:protecting-edf}).
\item \textbf{Formats read by an add-on}, such as the \file{.Stp} files of BedMasterEx, are read
  through the recording's adapter (\cref{chap:addons}), which answers \code{identify}.
\end{itemize}
A file that names several patients is flagged. The identifiers are kept in the table
\code{recording\_identifiers} only, and are shown only to users with the PHI permission who are
out of Privacy Mode: in the Encounters list's \ui{In the files}, in the encounter panel, in
\ui{Import} $\rightarrow$ \ui{Recordings}, and through the searches \code{file-name:},
\code{file-id:} and \code{file-dob:}, so that a reviewer can check the databases against them.
They never reach job summaries, the viewers, the actions' job descriptions, the results folders
or the logs.

\begin{phibox}
What a recording's file says about its patient\phimark{} is protected health information, like
the identifiers of the databases: it is kept in the catalog's database and nowhere else, and a
link that a file contradicts is listed in the Review tab for a person to decide.
\end{phibox}
